\documentclass[a4paper,12pt]{article}
\usepackage{amsmath,amssymb,amsthm}
\usepackage[margin=1in]{geometry}

\newtheorem{theorem}{Theorem}
\newtheorem{lemma}[theorem]{Lemma}
\newtheorem{corollary}[theorem]{Corollary}

\title{Problem 354: Distances in a Bee's Honeycomb}
\author{Project Euler}
\date{}

\begin{document}
\maketitle

\section{Problem Statement}

On a hexagonal lattice, define the Eisenstein norm of $z = a + b\omega$ (where $\omega = e^{2\pi i/3}$) as $N(z) = a^2 - ab + b^2$. Let $C(n)$ be the number of Eisenstein integers of norm exactly~$n$. Find the number of values $n \le 10^{12}$ with $C(n) = 450$.

\section{Mathematical Foundation}

\begin{theorem}[Representation count]
The number of Eisenstein integers of norm~$n$ is
\[
C(n) = 6\sum_{d \mid n} \chi(d),
\]
where $\chi$ is the non-principal Dirichlet character modulo~3.
\end{theorem}

\begin{proof}
The Eisenstein integers $\mathbb{Z}[\omega]$ form a unique factorization domain with 6~units. The norm is multiplicative, and the representation count factors as $6\sum_{d \mid n}\chi(d)$ by standard theory of quadratic forms over $\mathbb{Z}[\omega]$ (see Ireland--Rosen, Ch.~9).
\end{proof}

\begin{lemma}[Multiplicative formula]
Write $n = 3^a \prod_i p_i^{a_i} \prod_j q_j^{b_j}$ with $p_i \equiv 1 \pmod{3}$ and $q_j \equiv 2 \pmod{3}$. Then
\[
\sum_{d \mid n}\chi(d) = \begin{cases}
\displaystyle\prod_i (a_i + 1) & \text{if all } b_j \text{ are even,} \\
0 & \text{otherwise.}
\end{cases}
\]
\end{lemma}

\begin{proof}
Since $\chi$ is completely multiplicative, $f(n) = \sum_{d \mid n}\chi(d)$ is multiplicative. Evaluating on prime powers:
\begin{itemize}
    \item $f(3^a) = 1$ (since $\chi(3) = 0$, only $d = 1$ contributes).
    \item $f(p^a) = a + 1$ for $p \equiv 1\pmod{3}$ (since $\chi(p) = 1$).
    \item $f(q^b) = \frac{1+(-1)^b}{2}$ for $q \equiv 2\pmod{3}$ (since $\chi(q) = -1$).
\end{itemize}
Multiplying over all prime-power factors gives the result.
\end{proof}

\begin{corollary}
$C(n) = 450$ if and only if all exponents of primes $\equiv 2\pmod{3}$ are even and $\prod_i(a_i+1) = 75$.
\end{corollary}

\begin{proof}
$C(n) = 450$ requires $6\prod_i(a_i+1) = 450$, i.e., $\prod_i(a_i+1) = 75$, together with the even-exponent condition from the lemma.
\end{proof}

\begin{lemma}[Factorizations of 75]
The unordered factorizations of $75 = 3 \times 5^2$ into factors $\ge 2$ are:
\[
\{75\},\quad \{25,3\},\quad \{15,5\},\quad \{5,5,3\}.
\]
The corresponding exponent tuples $(a_i)$ are $(74)$, $(24,2)$, $(14,4)$, $(4,4,2)$.
\end{lemma}

\begin{proof}
Exhaustive enumeration.
\end{proof}

\section{Algorithm}

\begin{enumerate}
    \item Generate primes up to $\sqrt{10^{12}} = 10^6$; classify as $\equiv 1$ or $\equiv 2\pmod{3}$.
    \item For each exponent tuple, enumerate valid~$n \le 10^{12}$ by backtracking:
    \begin{itemize}
        \item Assign primes $\equiv 1\pmod{3}$ (in increasing order) to exponent slots.
        \item Multiply by arbitrary powers of~3 and even powers of primes $\equiv 2\pmod{3}$.
        \item Prune when the partial product exceeds $10^{12}$.
    \end{itemize}
    \item Sum the counts over all tuples.
\end{enumerate}

\section{Complexity Analysis}

\begin{itemize}
    \item \textbf{Time:} Dominated by the backtracking enumeration; aggressive pruning keeps runtime sub-second since only 58{,}065 valid~$n$ exist.
    \item \textbf{Space:} $O(\pi(10^6))$ for the prime list, plus $O(k)$ recursion depth.
\end{itemize}

\section{Answer}

\[
\boxed{58065134}
\]

\end{document}
